\documentclass[a4paper,fontset=none]{ctexart}

\usepackage{indentfirst}
\setlength{\parindent}{0em}
\usepackage{autobreak}
\allowdisplaybreaks
\usepackage{parskip}
\usepackage{geometry}
\geometry{top=3cm,bottom=3cm,left=2cm,right=2cm}
\usepackage{anyfontsize}

\usepackage{fontspec}
\setmainfont{STIX Two Text}
\usepackage{unicode-math}
\unimathsetup{mathrm=sym,mathbf=sym,mathsf=sym,bold-style=ISO}
\setmathfont{STIX Two Math}
\usepackage{physics2}
\usephysicsmodule{ab,op.legacy}
\usepackage{fixdif,derivative}
\newcommand{\um}{\upmu\mathrm{m}}
\newcommand{\ang}{\text{\r{A}}}
\AtBeginDocument{
    \let\uglyepsilon\epsilon
    \renewcommand\epsilon{\varepsilon}
}

\setCJKmainfont{Source Han Serif CN}[BoldFont=Source Han Sans CN Medium,ItalicFont=FZKai-Z03]

\title{\textbf{星际介质天文学作业}}
\author{张植竣}
\date{2024年5月24日}

\begin{document}

\maketitle

\begin{enumerate}
    \item[39.1]
    假设把能量$E$注入密度为$\rho_0$的均匀密度零温介质，其中一半的能量变为介质的动能，则有：
    \[
        \frac{1}{2}E = \frac{1}{2}Mv_s^2
    \]
    又因为：
    \[
        M = \frac{4\pi}{3}R_s^3,\quad R_s \propto t^\eta,\quad v_s = \odv{R_s}{t} = \eta\frac{R_s}{t}
    \]
    则有：
    \[
        E = Mv_s^2 = \frac{4\pi}{3}R_s^3 \ab(\eta\frac{R_s}{t})^2
    \]
    解得$\eta = 2/5$，即：
    \[
        R_s = A'E^{1/5}\rho_0^{-1/5}t^{2/5}
    \]
    其中：
    \[
        A' = \ab(\frac{3}{4\pi\eta^2})^{1/5} = \ab(\frac{75}{16\pi})^{1/5} = 1.08332
    \]
    与精确值$A = 1.15167$相比非常接近，相差约$6\%$。

    \bigskip
    \item[39.2 (a)]
    假设：
    \[
        R_s = A E^\alpha \rho^\beta t^\eta
    \]
    由量纲分析可知$\alpha = 1/5$、$\beta = -1/5$、$\eta = 2/5$，即：
    \[
        R_s = A E^{1/5} \rho^{-1/5} t^{2/5}
    \]
    当环境密度$\rho = \rho_0(r/r_0)^\delta$，能量$E = E_0(t/t_0)^\epsilon$时，上式应当仍成立，即：
    \begin{align*}
        R_s
        &= A \ab[E_0 \ab(\frac{t}{t_0})^\epsilon]^\alpha \ab[\rho_0 \ab(\frac{r}{r_0})^\delta]^\beta t^\eta \\
        &= A \ab[E_0 \ab(\frac{t}{t_0})^\epsilon]^{1/5} \ab[\rho_0 \ab(\frac{r}{r_0})^\delta]^{-1/5} t^{2/5} \\
        &= A E_0^{1/5} t^{\epsilon/5} t_0^{-\epsilon/5} \rho_0^{-1/5} r^{-\delta/5} r_0^{\delta/5} t^{2/5}
    \end{align*}
    注意这里的$r$即为$R_s$，则有：
    \[
        R_s^{1+\delta/5} = A E_0^{1/5} t^{\epsilon/5} t_0^{-\epsilon/5} \rho_0^{-1/5} r_0^{\delta/5} t^{2/5}
    \]
    整理得：
    \[
        R_s = \ab(A E_0 t_0^{-\epsilon} \rho_0^{-1} r_0^\delta)^{\tfrac{1}{5+\delta}} \cdot t^{\tfrac{2+\epsilon}{5+\delta}} \propto t^{\tfrac{2+\epsilon}{5+\delta}}
    \]
    与$R \propto t^\gamma$对比可知：
    \[
        \gamma = \frac{2+\epsilon}{5+\delta}
    \]

    \medskip
    \item[39.2 (b)]
    当激波的半径$R_s \propto t^\gamma$时，设$R_s = Ct^\gamma$，则激波的速度：
    \[
        v_s = \odv{R_s}{t} = C\gamma t^{\gamma - 1}
    \]
    则激波温度为：
    \[
        T_s = \frac{3}{16} \frac{\mu v_s^2}{k} = \frac{3\mu C^2}{16k} \gamma^2 t^{2(\gamma - 1)}
    \]
    其正比于$t^{2(\gamma - 1)}$。

    \medskip
    \item[39.2 (c)]
    由题意，$\delta = -2$，$\epsilon = 0$，则：
    \[
        \gamma = \frac{2+\epsilon}{5+\delta} = \gamma = \frac{2+0}{5-2} = \frac{2}{3}
    \]
\end{enumerate}

\end{document}